\documentclass[12pt]{article}
\usepackage[utf8]{inputenc}
\usepackage[T1]{fontenc}
\usepackage{geometry}
\geometry{margin=1in}
\usepackage{hyperref}
\usepackage{parskip}

\begin{document}

\begin{flushright}
March 26, 2026
\end{flushright}

\vspace{0.5cm}

Dear Editor-in-Chief of \textit{Technovation},

\vspace{0.3cm}

We are pleased to submit our manuscript entitled ``\textbf{Technology Adoption in Dual Economies: How Informality and Non-Wage Labor Costs Shape Automation Dynamics --- Evidence from Colombia}'' for consideration as a research article in \textit{Technovation}.

We present the first quasi-experimental evidence that labor cost shocks \emph{reduce} automation investment in dual economies, contradicting the standard theoretical prediction that higher labor costs accelerate technology adoption. Using a difference-in-differences design exploiting minimum wage shocks of +16\% (2023) and +12\% (2024) under Colombia's Petro administration, we find that firms most exposed to these shocks reduced automation investment per worker by 16.7\% ($\beta = -0.167$, $p = 0.045$). Drawing on 66,775 firm-year observations from Colombia's Annual Manufacturing Survey (2016--2024) and 364,998 individual observations from household labor force data, we document the ``dual economy trap'': a mechanism by which cash constraints, informality as an escape valve, and regulatory uncertainty cause wage increases to consume firm margins rather than redirecting capital toward labor-replacing technologies.

We believe this paper is well suited to \textit{Technovation} for several reasons. The manuscript directly addresses technology adoption dynamics---a core concern of the journal---by identifying institutional constraints on technology diffusion in developing countries. Our findings challenge the universality of OECD-calibrated automation models and demonstrate that innovation barriers in economies with large informal sectors require fundamentally different analytical frameworks. The paper speaks to the journal's interest in the interaction between technology, institutions, and economic development.

The paper makes three principal contributions. First, it provides the first causal evidence on the effect of labor cost shocks on automation investment in a dual economy, exploiting exogenous variation from minimum wage policy that no prior study has identified in this setting. Second, we articulate and formalize the ``dual economy trap'' as a novel theoretical mechanism explaining why the cost--automation nexus reverses in cash-constrained economies with large informal sectors. Third, principal component analysis validates a sectoral vulnerability index showing that automation potential and formality---not labor cost alone---drive vulnerability rankings, with direct implications for policy prioritization in developing economies.

This manuscript is original work that has not been published elsewhere and is not currently under review at any other journal. All authors have read and approved the final version of the manuscript.

We look forward to your consideration of this work.

\vspace{0.5cm}

Sincerely,

\vspace{0.5cm}

Cristian Espinal Maya\\
Department of Economics, Universidad EAFIT\\
Carrera 49 No. 7 Sur-50, Medell\'in 050022, Colombia\\
\href{mailto:cjespinalm@eafit.edu.co}{cjespinalm@eafit.edu.co}

\end{document}
